% Author: Grayson Orr
% Course: IN607: Introduction Application Development Concepts

\documentclass{article}
\author{}

\usepackage{graphicx}
\usepackage{wrapfig}
\usepackage{enumerate}
\usepackage{hyperref}
\usepackage[margin = 2.25cm]{geometry}
\usepackage[table]{xcolor}
\hypersetup{
  colorlinks = true,
  urlcolor = blue
}
\setlength\parindent{0pt}

\begin{document}

\begin{figure}
	\includegraphics[width=50mm]{./img/logo.png}
\end{figure}

\title{Course Directive\\IN607: Introduction Application Development Concepts\\Semester One, 2021}
\date{}
\maketitle

\section*{Course Information}
\begin{tabular}{ll}
	Credits:            & 15 Credits                             \\
	Prerequisite:       & IN511: Programming 2                   \\
	Recommended:        & IN512: Fundamentals of Web Development \\
	Stream A Timetable: & Add times here \& Add times here       
\end{tabular} 

\section*{Lecturers}
\begin{tabular}{ll}
	Name:   & Adon Moskal (Principal Lecturer) \\
	Office: & D205b                            \\
	Email:  & adon.moskal@op.ac.nz             \\
\end{tabular}
\begin{tabular}{l}
	Grayson Orr (Lecturer) \\
	D311                   \\
	grayson.orr@op.ac.nz   \\
\end{tabular}

\section*{Course Dates}
\begin{tabular}{ll}
	Term 1:             & 22 February - 16 April (8 weeks) \\
	Mid Semester Break: & 19 April - 30 April (2 weeks)    \\
	Term 2:             & 3 May - 25 February (8 weeks)    \\                       
\end{tabular}

\section*{Aims}
To introduce the concepts of application development including algorithms, data structures \& design patterns that are required to use a simple, industry-relevant development framework.

\section*{Learning Outcomes}
At the successful completion of this course, student will be able to:
\begin{enumerate}
	\item Design \& build usable, secure \& attractive applications with dynamic database
	      functionality following an appropriate software development methodology.
\end{enumerate}

\section*{Resources}

\subsection*{Software}
This paper will be taught using \textbf{Microsoft Visual Studio 2019} \& \textbf{Microsoft Visual Studio Code}. An installer for \textbf{Microsoft Visual Studio 2019} \&  \textbf{Microsoft Visual Studio Code} is available. See \href{https://visualstudio.microsoft.com/downloads}{https://visualstudio.microsoft.com/downloads} \& \href{https://code.visualstudio.com/Download}{https://code.visualstudio.com/Download}. Please refer any problems with downloads or installers to \textbf{Rob Broadley} in \textbf{D205a}.

\subsection*{Readings}
There is no textbook for the course.

\section*{Provisional Schedule}

\renewcommand{\arraystretch}{1.5}
\begin{tabular}{|c|c|c|c|}
	\hline
	\textbf{Week} & \textbf{Date}     & \multicolumn{2}{c|}{\textbf{Session}}        \\ \hline
	\small 1      & \small 22-02-2020 & \multicolumn{2}{c|}{\small ASP.NET Core 5.0} \\ \hline
	\small 2      & \small 01-03-2020 & \multicolumn{2}{c|}{\small ASP.NET Core 5.0} \\ \hline
	\small 3      & \small 08-03-2020 & \multicolumn{2}{c|}{\small MongoDB}          \\ \hline
	\small 4      & \small 15-03-2020 & \multicolumn{2}{c|}{\small JSON Web Tokens}  \\ \hline
	\small 5      & \small 22-03-2020 & \multicolumn{2}{c|}{\small NUnit}            \\ \hline
	\small 6      & \small 29-03-2020 & \multicolumn{2}{c|}{\small React}            \\ \hline
	\small 7      & \small 05-04-2020 & \multicolumn{2}{c|}{\small React}            \\ \hline
	\small 8      & \small 12-04-2020 & \multicolumn{2}{c|}{\small Cypress.IO}       \\ \hline
	\rowcolor{yellow} \multicolumn{4}{|c|}{\small Mid Term Break}  \\ \hline
	\small 9      & \small 03-05-2020 & \multicolumn{2}{c|}{\small Azure}            \\ \hline
	\small 10     & \small 10-05-2020 & \multicolumn{2}{c|}{\small Project Work}     \\ \hline
	\small 11     & \small 17-05-2020 & \multicolumn{2}{c|}{\small Project Work}     \\ \hline
	\small 12     & \small 24-05-2020 & \multicolumn{2}{c|}{\small Project Work}     \\ \hline
	\small 13     & \small 31-05-2020 & \multicolumn{2}{c|}{\small Project Work}     \\ \hline
	\small 14     & \small 07-06-2020 & \multicolumn{2}{c|}{\small Project Work}     \\ \hline
	\small 15     & \small 14-06-2020 & \multicolumn{2}{c|}{\small Project Work}     \\ \hline
	\small 16     & \small 21-06-2020 & \multicolumn{2}{c|}{\small Project Work}     \\ \hline
\end{tabular}

\section*{Assessments}
\renewcommand{\arraystretch}{1.5}
\begin{tabular}{|c|c|c|c|}
	\hline
	\textbf{Assessment}                                & \textbf{Weighting} & \textbf{Due Date} & \textbf{Learning Outcomes} \\ \hline 
	\small Project                       & \small 100\%        & \small 23-06-2020 & \small 1                   \\ \hline
\end{tabular}

\section*{Course Requirements \& Expectations}

\subsection*{Learning Hours}
This course requires 150 hours of learning. This time includes 64 hours of timetabled class time, \& 86 hours of self-directed reading, preparation \& completion of assessment work.

\subsection*{Criteria for Passing}
To pass this paper, you must achieve a cumulative mark of 50\%. Resits \& reassessments are not applicable.

\subsection*{Attendance}
\begin{itemize}
	\item You are expected to attend all classes, i.e, lectures \& labs.
	\item If you miss a class, you will need to catch up in your own self-directed time.
	\item If you cannot attend a class for any reason, please contact your lecturers (Adon Moskal \& Grayson Orr).
\end{itemize}

\subsection*{Communication}
Microsoft Outlook \& Teams are the official communication channels. It is your responsibility to regularly check Microsoft Outlook/Teams \& \href{https://github.com/otago-polytechnic-bit-courses/intro-app-dev-concepts}{GitHub} for important course-related material, including changes to class scheduling or assessment details. Not checking will not be accepted as an excuse.

\subsection*{Snow Days/Polytechnic Closure}
In the event the Polytechnic is closed or has a delayed opening because of snow or bad weather, you should not attempt to attend class if it is unsafe to do so. It is possible that your lecturer will not be able to attend either, so classes will not physically be meeting. However, this does not become a holiday. Rather, the material will be made available on \href{https://github.com/otago-polytechnic-bit-courses/intro-app-dev-concepts}{GitHub} for classes affected by the closure. You are responsible for any material presented in this manner. Information about closure will be posted on the \href{https://www.facebook.com/OtagoPoly}{Otago Polytechnic Facebook} page.

\subsection*{Group Work \& Originality}
Students in the Bachelor of Information Technology degree are expected to hand in original work. Students are encouraged to discuss assessments with their fellow students, however, all assessments are to be completed as individual work unless group work is explicitly required (i.e. if it doesn’t say it is group work then it is not group work – even if a group consultation was involved). Failure to submit your own original work will be treated as plagiarism.

\subsection*{Plagiarism}
Plagiarism is submitting someone elses work as your own. Plagiarism offences are taken seriously \& an assessment that has been plagiarised may be awarded zero. A definition of plagiarism is available \href{https://studentservices.op.ac.nz/learning-support/academic-integrity/}{here}.

\subsection*{Referencing}
Appropriate referencing is required for all work. Referencing standards will be specified by your lecturers (Adon Moskal \& Grayson Orr).

\subsection*{Submission Requirements}
All assessments are to be submitted by the time, date \& method given when the assessment is issued. Failure to meet all requirements may result in a penalty of up to 10\% per day (including weekends).

\subsection*{Extensions}
Extensions are only available for unusual circumstances. These must be applied for \& approved, before the assessment due date. 

\subsection*{Impairment}
In case of sickness contact your lecturers (Adon Moskal \& Grayson Orr) or BIT Team Leader (Michael Holtz) as soon as possible, preferably before the assessment or exam is due. The policy regarding the granting of a mark that considers impaired performance requires a medical certificate \& a medical practitioner’s signature on a form. You may refer to the guide on impaired performance on the student handbook.

\subsection*{Appeals}
If you are concerned about any aspect of your assessment, please approach lecturers (Adon Moskal \& Grayson Orr) in the first instance. We support an open-door policy \& aim to resolve issues promptly. Further support is available from the BIT Team Leader (Michael Holtz) \& Head of College (Richard Nyhof). Otago Polytechnic has a formal process for academic appeals if necessary.

\subsection*{Other Documents}
Regulatory documents relating to this course can be found on the Otago Polytechnic website.

\end{document}